% Unit 117 terjemahan Bahasa Indonesia dari chapter8.tex baris 1704--1849.
% Sumber: Methods of Algebra, Volume 2, commit
% 9a5803ff77dd3257484cb177f851a73770a59dd3 (CC BY 4.0).
% stable-unit-id: o014.aljabr2.chapter8.closed-structure
% Cakupan: Bagian 8.9 lengkap, Struktur Tertutup, hingga tepat sebelum
% Bagian 8.10.
% Perbaikan matematis-redaksional: pembukaan membedakan ketertutupan
% Kartesius sSet dari ketertutupan monoidal sAb; arah f^* pada baris sumber
% 1710 dibalik agar sesuai dengan prakomposisi dan baris 1717--1722.

% segment-id: o014.aljabr2.chapter8.closed-structure.h001
\section{Struktur Tertutup}\label{sec:closed-simplicial}
\hypertarget{o014.aljabr2.chapter8.closed-structure}{ }
% segment-id: o014.aljabr2.chapter8.closed-structure.p001
Tujuan bagian ini adalah menjelaskan struktur tertutup pada $\cate{sSet}$
dan $\cate{sAb}$. Untuk $\cate{sSet}$, struktur ini bersifat tertutup
Kartesius dalam pengertian Definisi~\ref{def:Cartesian-closed}; untuk
$\cate{sAb}$, struktur monoidal dari hasil kali tensorlah yang tertutup.
Dengan kata lain, kita hendak mengangkat himpunan $\Hom$ menjadi objek
simplisial yang sesuai. Pertama-tama kita membahas kategori himpunan
simplisial $\cate{sSet}$, kemudian versinya untuk $\cate{sAb}$ dan
perbandingannya, melalui korespondensi Dold--Kan, dengan kompleks $\Hom$.
Langkah terakhir memerlukan hasil dari \S\ref{sec:bisimplicial}.

% segment-id: o014.aljabr2.chapter8.closed-structure.d001
\begin{definition}\label{def:internal-Hom-sSet}
	Untuk $X,Y\in\Obj(\cate{sSet})$, definisikan
	$\iHom(X,Y)=\iHom_{\cate{sSet}}(X,Y)\in\Obj(\cate{sSet})$ sebagai berikut:
	% segment-id: o014.aljabr2.chapter8.closed-structure.q001
	\[
	\iHom(X,Y)_n:=\Hom_{\cate{sSet}}(X\times\Delta^n,Y),
	\qquad n\in\Z_{\geq0}.
	\]
	Untuk setiap morfisme $f:[m]\to[n]$ dalam $\simpDelta$, definisikan
	$f^*:\iHom(X,Y)_n\to\iHom(X,Y)_m$ sebagai penarikan balik, yakni
	prakomposisi, sepanjang
	% segment-id: o014.aljabr2.chapter8.closed-structure.q002
	\[
	\identity_X\times f:X\times\Delta^m\to X\times\Delta^n.
	\]
	Definisikan pula morfisme evaluasi
	% segment-id: o014.aljabr2.chapter8.closed-structure.q003
	\[
	\mathrm{ev}_{X,Y}:\iHom(X,Y)\times X\to Y
	\]
	sebagai berikut: citra
	$(\varphi,x)\in\Hom_{\cate{sSet}}(X\times\Delta^n,Y)\times X_n$ ialah
	$\mathrm{ev}_{X,Y,n}(\varphi,x):=\varphi(x,\identity_{[n]})\in Y_n$.
	Ingat bahwa $(\Delta^n)_n=\End_{\simpDelta}([n])$.
\end{definition}

% segment-id: o014.aljabr2.chapter8.closed-structure.p002
Kita harus memeriksa bahwa $\mathrm{ev}_{X,Y}$ memang morfisme dalam
$\cate{sSet}$. Pemeriksaan ini langsung: untuk $f:[m]\to[n]$ dan
$(\varphi,x)\in\iHom(X,Y)_n\times X_n$, berlaku
% segment-id: o014.aljabr2.chapter8.closed-structure.q004
\begin{align*}
	\mathrm{ev}_{X,Y,m}\bigl(f^*(\varphi),f^*(x)\bigr)
	&=((\identity_X\times f)^*\varphi)
	  \bigl(f^*(x),\identity_{[m]}\bigr)\\
	&=\varphi\bigl(
	  \underbracket{f^*(x),\;f}_{\in X_m\times(\Delta^n)_m}\bigr)
	 =\varphi\bigl(f^*(x),f^*\identity_{[n]}\bigr)\\
	&=f^*\bigl(\varphi(x,\identity_{[n]})\bigr)
	 =f^*\bigl(\mathrm{ev}_{X,Y,n}(\varphi,x)\bigr).
\end{align*}

% segment-id: o014.aljabr2.chapter8.closed-structure.p003
Dengan memvariasikan $X,Y$, konstruksi ini memberikan funktor
$\iHom:\cate{sSet}^{\opp}\times\cate{sSet}\to\cate{sSet}$, sedangkan
$\mathrm{ev}$ alami terhadap $X$ dan $Y$.

% segment-id: o014.aljabr2.chapter8.closed-structure.pr001
\begin{proposition}\label{prop:sSet-closed}
	Untuk setiap $X,Y,Z\in\Obj(\cate{sSet})$, terdapat bijeksi kanonik
	% segment-id: o014.aljabr2.chapter8.closed-structure.g001
	\[\begin{tikzcd}
		\Hom_{\cate{sSet}}\bigl(X,\iHom(Y,Z)\bigr)
		\arrow[r,"1:1"]
		&\Hom_{\cate{sSet}}(X\times Y,Z).
	\end{tikzcd}\]
	Bijeksi ini memetakan $g:X\to\iHom(Y,Z)$ ke komposisi
	% segment-id: o014.aljabr2.chapter8.closed-structure.q005
	\[
	X\times Y\xrightarrow{g\times\identity_Y}
	\iHom(Y,Z)\times Y\xrightarrow{\mathrm{ev}_{Y,Z}}Z.
	\]
	Sebaliknya, bijeksi tersebut memetakan $h:X\times Y\to Z$ ke morfisme
	berikut. Untuk $n\in\Z_{\geq0}$ dan $x\in X_n$, misalkan
	$\iota_x:\Delta^n\to X$ morfisme yang bersesuaian dengan $x$. Citra $x$
	ialah komposisi
	% segment-id: o014.aljabr2.chapter8.closed-structure.q006
	\[
	Y\times\Delta^n\xrightarrow{\identity_Y\times\iota_x}Y\times X
	\xrightarrow[\sim]{\text{pertukaran}}X\times Y\xrightarrow{h}Z.
	\]
\end{proposition}
% segment-id: o014.aljabr2.chapter8.closed-structure.pf001
\begin{proof}
	Pemeriksaan rutin.
\end{proof}

% segment-id: o014.aljabr2.chapter8.closed-structure.c001
\begin{corollary}
	Bijeksi dalam Proposisi~\ref{prop:sSet-closed} menjadikan $\cate{sSet}$
	kategori tertutup Kartesius dalam pengertian
	Definisi~\ref{def:Cartesian-closed}, dengan bifunktor $\iHom$.
\end{corollary}
% segment-id: o014.aljabr2.chapter8.closed-structure.pf002
\begin{proof}
	Ingatlah bahwa hasil kali dalam $\cate{sSet}$ tidak lain adalah hasil kali himpunan
	simplisial $X\times Y$, dan sebagainya. Tinggal membandingkan
	\eqref{eqn:closed-monoidal-adj0} dengan bijeksi kanonik dalam
	Proposisi~\ref{prop:sSet-closed}.
\end{proof}

% segment-id: o014.aljabr2.chapter8.closed-structure.p004
Teori umum kategori monoidal tertutup (lihat
\S\ref{sec:closed-monoidal}) memberikan morfisme kanonik
$\iHom(X,Y)\times X\to Y$: morfisme ini adalah citra $\identity$ di bawah
% segment-id: o014.aljabr2.chapter8.closed-structure.q007
\[
\End_{\cate{sSet}}\bigl(\iHom(X,Y)\bigr)\xrightarrow{1:1}
\Hom_{\cate{sSet}}\bigl(\iHom(X,Y)\times X,Y\bigr).
\]
Dengan menguraikan definisi, terlihat bahwa morfisme tersebut tepat sama
dengan $\mathrm{ev}_{X,Y}$ yang didefinisikan sebelumnya.

% segment-id: o014.aljabr2.chapter8.closed-structure.eg001
\begin{example}
	Proposisi~\ref{prop:nerve-ff} menyatakan bahwa funktor saraf $\mathrm{N}$
	menanamkan $\cate{Cat}$ secara setia penuh ke dalam $\cate{sSet}$.
	Bagaimana $\iHom$ tercermin pada tingkat kategori? Jawabannya sederhana:
	untuk setiap kategori kecil $\mathcal{C}$ dan $\mathcal{D}$, terdapat
	isomorfisme kanonik
	% segment-id: o014.aljabr2.chapter8.closed-structure.q008
	\[
	\iHom(\mathrm{N}\mathcal{C},\mathrm{N}\mathcal{D})
	\simeq\mathrm{N}(\mathcal{D}^{\mathcal{C}}).
	\]
	Memang, $\mathrm{N}$ memetakan produk kategori ke produk himpunan
	simplisial, dan $\mathrm{N}([n])=\Delta^n$. Karena $\cate{Cat}$ juga
	tertutup Kartesius, diperoleh
	% segment-id: o014.aljabr2.chapter8.closed-structure.q009
	\begin{multline*}
		\Hom_{\cate{sSet}}(\mathrm{N}\mathcal{C}\times\Delta^n,
		\mathrm{N}\mathcal{D})
		=\Hom_{\cate{sSet}}(\mathrm{N}(\mathcal{C}\times[n]),
		\mathrm{N}\mathcal{D})\\
		\simeq\Hom_{\cate{Cat}}(\mathcal{C}\times[n],\mathcal{D})
		\simeq\Hom_{\cate{Cat}}([n],\mathcal{D}^{\mathcal{C}})
		=\mathrm{N}(\mathcal{D}^{\mathcal{C}})_n.
	\end{multline*}
	Mudah diperiksa bahwa untuk setiap morfisme $f:[m]\to[n]$ dalam
	$\simpDelta$, penarikan balik pada kedua ruas saling bersesuaian.
\end{example}

% segment-id: o014.aljabr2.chapter8.closed-structure.p005
Sekarang tinjau $\cate{sAb}$. Ingat bahwa kategori ini mempunyai struktur
monoidal simetris yang berasal dari $\otimes:=\otimes_{\Z}$, yakni dengan
mengambil hasil kali tensor suku demi suku. Untuk
$X\in\Obj(\cate{sAb})$ dan $K\in\Obj(\cate{sSet})$, definisikan
% segment-id: o014.aljabr2.chapter8.closed-structure.q010
\[
X\otimes K:=X\otimes\Z K\in\Obj(\cate{sAb}).
\]
Khususnya, $X\otimes\Delta^n$ terdefinisi. Konstruksi ini memberikan
bifunktor $\cate{sAb}\times\cate{sSet}\to\cate{sAb}$.

% segment-id: o014.aljabr2.chapter8.closed-structure.d002
\begin{definition}
	Untuk $X,Y\in\Obj(\cate{sAb})$, definisikan
	$\iHom(X,Y)=\iHom_{\cate{sAb}}(X,Y)\in\Obj(\cate{sAb})$ melalui
	$\iHom(X,Y)_n:=\Hom_{\cate{sAb}}(X\otimes\Delta^n,Y)$.
	Morfisme-morfisme penarikan balik didefinisikan seperti dalam
	Definisi~\ref{def:internal-Hom-sSet}.
\end{definition}

% segment-id: o014.aljabr2.chapter8.closed-structure.pr002
\begin{proposition}\label{prop:sAb-closed}
	Untuk setiap $X,Y,Z\in\Obj(\cate{sAb})$, terdapat diagram komutatif kanonik
	% segment-id: o014.aljabr2.chapter8.closed-structure.g002
	\[\begin{tikzcd}
		\Hom_{\cate{sSet}}\bigl(X,\iHom_{\cate{sSet}}(Y,Z)\bigr)
		\arrow[r,"1:1"]
		&\Hom_{\cate{sSet}}(X\times Y,Z)\\
		\Hom_{\cate{sAb}}\bigl(X,\iHom_{\cate{sAb}}(Y,Z)\bigr)
		\arrow[r,"1:1"']
		\arrow[phantom,u,"\subset" description,sloped]
		&\Hom_{\cate{sAb}}(X\otimes Y,Z)\arrow[hookrightarrow,u]
	\end{tikzcd}\]
\end{proposition}
% segment-id: o014.aljabr2.chapter8.closed-structure.pf003
\begin{proof}
	Baris pertama diagram adalah Proposisi~\ref{prop:sSet-closed}. Panah
	vertikal di sebelah kanan menanamkan
	$\Hom_{\cate{sAb}}(X\otimes Y,Z)$ sebagai bagian bilinear dari
	$\Hom_{\cate{sSet}}(X\times Y,Z)$. Citra penanaman di sebelah kiri terdiri
	atas semua $\varphi:X\to\iHom_{\cate{sSet}}(Y,Z)$ dengan keluarga
	morfisme
	% segment-id: o014.aljabr2.chapter8.closed-structure.q012
	\[
	\varphi_n:X_n\to\Hom_{\cate{sSet}}(Y\times\Delta^n,Z)
	\]
	memenuhi, untuk setiap $n\in\Z_{\geq0}$,
	% segment-id: o014.aljabr2.chapter8.closed-structure.l001
	\begin{itemize}
		\item $\varphi_n(x):Y\times\Delta^n\to Z$ aditif dalam variabel $Y$
		untuk setiap $x\in X_n$, atau dengan kata lain berfaktor melalui
		$\Hom_{\cate{sAb}}(Y\otimes\Delta^n,Z)$;
		\item peta hasil faktorisasi
		$\varphi_n:X_n\to\Hom_{\cate{sAb}}(Y\otimes\Delta^n,Z)$ bersifat
		aditif.
	\end{itemize}
	Dengan menguraikan definisi, langsung terlihat bahwa kedua deskripsi ini
	saling bersesuaian.
\end{proof}

% segment-id: o014.aljabr2.chapter8.closed-structure.p006
Sesungguhnya, $\cate{sAb}$ adalah kategori aditif, dan bijeksi
$\Hom_{\cate{sAb}}\bigl(X,\iHom_{\cate{sAb}}(Y,Z)\bigr)
\simeq\Hom_{\cate{sAb}}(X\otimes Y,Z)$ merupakan isomorfisme modul-$\Z$.

% segment-id: o014.aljabr2.chapter8.closed-structure.c002
\begin{corollary}
	Misalkan $X,Y\in\Obj(\cate{sAb})$ dan $K\in\Obj(\cate{sSet})$. Terdapat
	isomorfisme kanonik modul-$\Z$
	% segment-id: o014.aljabr2.chapter8.closed-structure.q013
	\begin{align*}
		\Hom_{\cate{sAb}}(X\otimes K,Y)
		&\simeq\Hom_{\cate{sSet}}
		\bigl(K,\iHom_{\cate{sAb}}(X,Y)\bigr)\\
		&\simeq\Hom_{\cate{sAb}}
		\bigl(X,\iHom_{\cate{sSet}}(K,Y)\bigr).
	\end{align*}
	Berdasarkan sifat universal, $\iHom_{\cate{sSet}}(K,Y)$ pada suku terakhir
	secara alami teridentifikasi dengan himpunan simplisial yang mendasari
	$\iHom_{\cate{sAb}}(\Z K,Y)$, sehingga memperoleh
	struktur objek $\cate{sAb}$.
\end{corollary}
% segment-id: o014.aljabr2.chapter8.closed-structure.pf004
\begin{proof}
	Proposisi~\ref{prop:sAb-closed} memberikan
	% segment-id: o014.aljabr2.chapter8.closed-structure.q014
	\begin{align*}
		\Hom_{\cate{sAb}}(X\otimes K,Y)
		&\simeq\Hom_{\cate{sAb}}(\Z K\otimes X,Y)\\
		&\simeq\Hom_{\cate{sAb}}
		\bigl(\Z K,\iHom_{\cate{sAb}}(X,Y)\bigr).
	\end{align*}
	Sifat adjungsi $\Z(\cdot)$ menunjukkan bahwa suku terakhir adalah
	$\Hom_{\cate{sSet}}(K,\iHom_{\cate{sAb}}(X,Y))$. Secara serupa,
	$\Hom_{\cate{sAb}}(X\otimes K,Y)\simeq
	\Hom_{\cate{sAb}}(X,\iHom_{\cate{sAb}}(\Z K,Y))$.
\end{proof}

% segment-id: o014.aljabr2.chapter8.closed-structure.c003
\begin{corollary}
	Bijeksi kanonik dalam Proposisi~\ref{prop:sAb-closed} menjadikan
	$(\cate{sAb},\otimes)$ kategori monoidal tertutup dalam pengertian
	Definisi~\ref{def:closed-monoidal-cat}, dengan bifunktor $\iHom$.
\end{corollary}

% segment-id: o014.aljabr2.chapter8.closed-structure.p007
Seperti pada himpunan simplisial, teori umum kategori monoidal tertutup
memberikan morfisme evaluasi kanonik
% segment-id: o014.aljabr2.chapter8.closed-structure.q015
\[
\mathrm{ev}_{X,Y}:\iHom_{\cate{sAb}}(X,Y)\otimes X\to Y.
\]
Morfisme ini mempunyai deskripsi konkret yang serupa dengan
Definisi~\ref{def:internal-Hom-sSet}; deskripsi tersebut diperoleh langsung
dengan menguraikan definisi.

% segment-id: o014.aljabr2.chapter8.closed-structure.p008
Sekarang kita kaji hubungan antara $\iHom_{\cate{sAb}}$ dan kompleks $\Hom$
dari sudut pandang korespondensi Dold--Kan
(Teorema~\ref{prop:Dold-Kan}). Keduanya masing-masing merupakan $\Hom$
internal dari kategori monoidal tertutup $\cate{sAb}$ dan
$\cate{Ch}_{\geq0}(\cate{Ab})$. Walaupun funktor kompleks rantai
ternormalisasi
$\mathrm{N}:\cate{sAb}\to\cate{Ch}_{\geq0}(\cate{Ab})$ merupakan
ekuivalensi, funktor itu bukan funktor monoidal. Karena itu, kita tidak dapat
langsung menyimpulkan adanya isomorfisme antara
$\mathrm{N}\iHom_{\cate{sAb}}(X,Y)$ dan
$\Hom^\bullet(\mathrm{N}X,\mathrm{N}Y)$. Meskipun demikian, hasil dari
\S\ref{sec:bisimplicial} tetap memberikan sepasang morfisme kanonik
% segment-id: o014.aljabr2.chapter8.closed-structure.g003
\begin{equation}\label{eqn:Hom-N}
\begin{tikzcd}
	\mathrm{N}\iHom_{\cate{sAb}}(X,Y)\arrow[shift left,r]
	&\Hom^\bullet(\mathrm{N}X,\mathrm{N}Y).\arrow[shift left,l]
\end{tikzcd}
\end{equation}
% segment-id: o014.aljabr2.chapter8.closed-structure.l002
\begin{itemize}
	\item Berdasarkan struktur monoidal tertutup
	$\cate{Ch}_{\geq0}(\cate{Ab})$, memberikan morfisme dari kiri ke kanan
	sama artinya dengan memberikan morfisme
	% segment-id: o014.aljabr2.chapter8.closed-structure.q016
	\[
	\mathrm{N}\iHom_{\cate{sAb}}(X,Y)\otimes\mathrm{N}X
	\longrightarrow\mathrm{N}Y.
	\]
	Ambil morfisme ini sebagai komposisi
	% segment-id: o014.aljabr2.chapter8.closed-structure.q017
	\[
	\mathrm{N}\iHom_{\cate{sAb}}(X,Y)\otimes\mathrm{N}X
	\xrightarrow{\mathrm{EZ}}
	\mathrm{N}\bigl(\iHom_{\cate{sAb}}(X,Y)\otimes X\bigr)
	\xrightarrow{\mathrm{N}(\mathrm{ev}_{X,Y})}\mathrm{N}Y.
	\]
	\item Terapkan funktor adjoin kiri $\Gamma$ dari $\mathrm{N}$, yang sekaligus
	kuasi-inversnya. Memberikan morfisme dari kanan ke kiri sama artinya dengan
	memberikan morfisme
	% segment-id: o014.aljabr2.chapter8.closed-structure.q018
	\[
	\Gamma\Hom^\bullet(\mathrm{N}X,\mathrm{N}Y)
	\longrightarrow\iHom_{\cate{sAb}}(X,Y).
	\]
	Berdasarkan struktur monoidal tertutup $\cate{sAb}$, ini juga sama artinya
	dengan memberikan morfisme
	$\Gamma\Hom^\bullet(\mathrm{N}X,\mathrm{N}Y)\otimes X\to Y$.
	Secara konkret, pilih morfisme tersebut sehingga citranya di bawah
	$\mathrm{N}$ sama dengan komposisi
	% segment-id: o014.aljabr2.chapter8.closed-structure.q019
	\begin{multline*}
		\mathrm{N}\bigl(\Gamma\Hom^\bullet(\mathrm{N}X,\mathrm{N}Y)
		\otimes X\bigr)
		\xrightarrow{\mathrm{AW}}
		\mathrm{N}\Gamma\Hom^\bullet(\mathrm{N}X,\mathrm{N}Y)
		\otimes\mathrm{N}X\\
		\simeq\Hom^\bullet(\mathrm{N}X,\mathrm{N}Y)\otimes\mathrm{N}X
		\xrightarrow{\text{evaluasi}}\mathrm{N}Y.
	\end{multline*}
\end{itemize}

% segment-id: o014.aljabr2.chapter8.closed-structure.p009
Inti konstruksi di atas adalah bahwa $\mathrm{N}$ merupakan funktor monoidal
longgar kanan sekaligus longgar kiri. Seandainya $\mathrm{N}$ merupakan
funktor monoidal, atau secara ekuivalen seandainya $\mathrm{EZ}$ dan
$\mathrm{AW}$ keduanya isomorfisme, argumen abstrak akan menunjukkan
bahwa kedua morfisme dalam \eqref{eqn:Hom-N} saling invers; tetapi keadaan
sebenarnya tidak demikian. Namun, Teorema Eilenberg--Zilber umum
\ref{prop:Eilenberg-Zilber} menyatakan bahwa $\mathrm{EZ}$ dan $\mathrm{AW}$
saling invers pada tingkat homologi. Dari sini diperoleh hubungan berikut
antara kompleks $\Hom$ dan $\iHom_{\cate{sAb}}$.

% segment-id: o014.aljabr2.chapter8.closed-structure.pr003
\begin{proposition}
	Untuk setiap $X,Y\in\Obj(\cate{sAb})$, pasangan morfisme dalam
	\eqref{eqn:Hom-N} saling invers pada tingkat homologi.
\end{proposition}

% segment-id: o014.aljabr2.chapter8.closed-structure.p010
Jika $\Bbbk$ gelanggang komutatif, hasil serupa tetap berlaku setelah
$\cate{Ab}$ diganti dengan kategori monoidal simetris $\Bbbk\dcate{Mod}$ dan
kita meninjau objek simplisial di dalamnya.

% segment-id: o014.aljabr2.chapter8.closed-structure.r001
\begin{remark}
	Definisi $\iHom$ dapat diperluas ke sebarang kategori aditif $\mathcal{A}$.
	Memang, dalam deskripsi $\iHom_{\cate{sAb}}(X,Y)_n$, berlaku
	% segment-id: o014.aljabr2.chapter8.closed-structure.q020
	\[
	(X\otimes\Z\Delta^n)_k
	=X_k\otimes\left(\bigoplus_{u:[k]\to[n]}\Z\right)
	\simeq\bigoplus_{u:[k]\to[n]}X_k.
	\]
	Suku terakhir hanya melibatkan jumlah langsung hingga dan tidak melibatkan
	unsur apa pun; pemetaan $d_i,s_j$, dan sebagainya juga mempunyai deskripsi
	sederhana dalam penyajian jumlah langsung ini. Karena itu, dapat
	didefinisikan bifunktor
	% segment-id: o014.aljabr2.chapter8.closed-structure.q021
	\[
	\iHom_{\mathcal{A}}:(\cate{s}\mathcal{A})^{\opp}
	\times\cate{s}\mathcal{A}\to\cate{s}\mathcal{A},
	\]
	yang aditif dalam setiap variabel.
\end{remark}
